\documentclass[11pt]{article}
\usepackage[margin=1in]{geometry}
\usepackage{amsmath,amssymb,amsthm}
\usepackage[T1]{fontenc}
\usepackage{lmodern}
\usepackage{xcolor}
\usepackage{booktabs}
\usepackage{longtable}
\usepackage{array}
\usepackage[hidelinks]{hyperref}
\usepackage{url}

\newtheorem{theorem}{Theorem}
\newtheorem{proposition}[theorem]{Proposition}
\newtheorem{lemma}[theorem]{Lemma}
\newtheorem{corollary}[theorem]{Corollary}
\theoremstyle{definition}
\newtheorem{definition}[theorem]{Definition}
\theoremstyle{remark}
\newtheorem{remark}[theorem]{Remark}

% Code-like text (not checked by a name gate).
\newcommand{\lean}[1]{\texttt{#1}}
% A declaration of one of our developments; every \decl{...} must be in that repository's gate.sh
% REQUIRED list (checked by the name gate).
\DeclareUrlCommand\decl{\urlstyle{tt}\def\UrlBreaks{\do\_\do\.}}
% A Mathlib declaration (pinned revision 0df444a3); not one of ours.
\DeclareUrlCommand\mdecl{\urlstyle{tt}\def\UrlBreaks{\do\_\do\.}}
\newcommand{\RR}{\mathbb{R}}
\newcommand{\CC}{\mathbb{C}}
\newcommand{\PH}[1]{\textcolor{red}{\textbf{[PLACEHOLDER:} #1\textbf{]}}}

% ---------------------------------------------------------------------------------------------
% PINS. Published 2026-10-14; every \src links to the file and line at this commit.
% ---------------------------------------------------------------------------------------------
% version-2 commits: clifford=927f7f82b3a1ae8342b1636b78bb9c04199c0d4b cliffordchain=b6649f48bc4eaf192f729efd26a617fa2aeddfef fivecircles=6040681ed19d180325f87162f690ebdad4a96961 lester=d307293ba23b92aece5f81b971704e6e11641bfa thebault=2edce1eecc7b97faefa8580bddaef955ab44bf73 kosnita=7eaae5300ce68cf3c0eb59fbc52d8456c4c7f819 blundon=50e1f17b92bd09d2fc515fb46133f2f84b46ff6c steinerlehmus=53b62463cb50388b51e7a703afeb08e7c8f1ebdc
\expandafter\def\csname commitof:clifford\endcsname{209abaa5e638f37aee8ec374e29b74341297cc46}
\expandafter\def\csname commitof:cliffordchain\endcsname{ba083b2158e1b96e2360b6e7d8fe04bd278ff5e2}
\expandafter\def\csname commitof:fivecircles\endcsname{f1fd9b7b2bd9f6b1e881425019f86fc49adba927}
\expandafter\def\csname commitof:lester\endcsname{9f94754987cbb9ddb8bf14bef295f2bfedb077df}
\expandafter\def\csname commitof:thebault\endcsname{c37f92e53c9ed790a2c90afdef95a7a35e6a7842}
\expandafter\def\csname commitof:kosnita\endcsname{702f0263a70136830aee7a9fd38ec7f4fc7d304d}
\expandafter\def\csname commitof:blundon\endcsname{118b4c9f3afd6d4120ea2c2d2e60c1e70d7dc8b8}
\expandafter\def\csname commitof:steinerlehmus\endcsname{a3a3d26bee6c68eae6d564228355524c46fe546d}
\newcommand{\commitof}[1]{\csname commitof:#1\endcsname}
\newcommand{\repo}[1]{https://github.com/jpbald93/#1}
\newcommand{\src}[4]{\ifcsname commitof:#1\endcsname
  \href{\repo{#1}/blob/\commitof{#1}/#2/#3.lean\#L#4}{\texttt{#3}:#4}\else\texttt{#3}:#4\fi}

\title{Circle configurations in classical plane geometry, formalized in Lean~4:\\
Clifford's chain, Miquel's five circles, Lester, Sawayama--Th\'ebault and Kosnita}
\author{Josh Bald\\ \small Independent researcher, Ontario, Canada\\ \small\texttt{jpbald93@gmail.com}\\
\small ORCID: \texttt{0009-0002-1317-6489}
\thanks{\copyright~2026 Josh Bald. Licensed under CC BY 4.0 (\url{https://creativecommons.org/licenses/by/4.0/}). The accompanying Lean code is licensed under the Apache License 2.0.}}
\date{14 October 2026; version~2, 22 October 2026 (adds Section~\ref{sec:l1415})}

\emergencystretch=3em
\begin{document}
\sloppy
\maketitle

\begin{abstract}
We give machine-checked proofs, in Lean~4 with Mathlib, of a family of classical theorems about
configurations of circles in the Euclidean plane. The centrepiece is the beginning of Clifford's
chain: for four circles through a common point in general position, the four circles through
triples of their second intersection points have a common point (the case $n=4$); for five such
circles the five points so obtained lie on one circle \emph{or one line} (the case $n=5$); and, under
a stated further non-degeneracy condition, for six circles the six circles so obtained have a common
point (the case $n=6$). We also prove Miquel's five-circles theorem for the pentagram (again with
the conclusion ``on a circle or on a line''), Lester's theorem that the circumcentre, the
nine-point centre and the two Fermat points of a scalene triangle are concyclic, the
Sawayama--Th\'ebault theorem on the centres of the two Th\'ebault circles and the incentre, and
Kosnita's concurrence theorem; and, as two short companions in the same style, Blundon's
inequalities and the Steiner--Lehmus theorem. Every degenerate configuration that a statement
excludes is named, and for several exclusions we exhibit configurations showing that the hypothesis
cannot simply be dropped. All proofs use only Lean's three standard axioms.
In the sources and search scope listed in Section~\ref{sec:related} we found no previous complete formal
proof of any of these theorems, with one partial exception: the algebraic core of Steiner--Lehmus
(over side lengths, with the bisector length given by its formula) has a Lean proof, and we found no
formal proof of the geometric theorem. The formal proofs we found closest to the circle theorems are of
the three-circle Miquel theorem, which is the base case of Clifford's chain. We also show, with an
explicit six-circle configuration, that the natural finite-plane formalization of the chain is false
for every $n\ge6$: its Clifford point lies at infinity. This concerns our predicate, whose general
position constrains only the initial pairs and triples; it is not a counterexample to the textbook
chain, which requires a genuine circle at every odd level, whereas here a level-$5$ object provably lies on
no circle. Version~2 adds twelve Mathlib-only restatements of the
main theorems, eleven with explicit witnesses (the Clifford $n=6$ witness and a restatement of the
general-$n$ classification remain open), and a comparison with textbook statements formalized
independently, which also proves, for the comparison's definitions, the existence of the Fermat points
and Th\'eb\-ault circles and Kosnita's theorem in the form ``concurrent or all parallel''.
\end{abstract}

\section{Introduction}

Many of the most attractive theorems of nineteenth-century plane geometry are \emph{incidence
theorems for circles}: some construction is carried out with circles, and at the end several points
turn out to lie on one circle, or several circles turn out to pass through one point. The simplest is
Miquel's pivot theorem \cite{miquel1838}: if a point is chosen on each side of a triangle, the three
circles through a vertex and the two chosen points on the adjacent sides pass through a common
point. It is formalized in Isabelle/HOL \cite{afpmiquel}, both as an incidence statement and with
existence of the common point; Rocq's HighSchoolGeometry \cite{hsgeometry} has the incidence form of
the related four-line (complete quadrilateral) Miquel theorem.

Clifford \cite{clifford1871} discovered that Miquel's theorem is the first step of an infinite chain.
Take $n$ circles in general position through a common point $M$. Any two meet again in a point;
any three determine a circle through their three such points; for four circles, the four circles so
obtained are concurrent; for five, the five points so obtained are concyclic; for six, the six circles
so obtained are concurrent; and so on, alternating between points and circles. Miquel's five-circles
theorem, the Sawayama--Th\'ebault theorem and Kosnita's theorem are further classical circle
configurations, treated in the standard texts \cite{johnson1929,pedoe1970} and, for Th\'ebault's
problem, by computer algebra \cite{chou1988}; Lester's theorem is more recent \cite{lester1997,yiu2010}.
In the sources we searched (Section~\ref{sec:related}) we found no complete formal proof of any of them
in a proof assistant.

\subsection{Contributions}
All results below are proved in Lean~4 \cite{lean4} with Mathlib \cite{mathlib2020}.
\begin{enumerate}
\item \textbf{Clifford's chain, $n=4,5,6$} (Theorems~\ref{thm:cliff4}--\ref{thm:cliff6}), in an explicit
complex-coordinate model, together with a second, independent development
(Section~\ref{sec:chain}) which defines an existential predicate for the chain at every $n$,
proves the cases $n=4,5$ in that general form ($n=5$ with no hypothesis beyond general position),
and proves that it fails for every $n\ge6$: an explicit six-circle configuration in general position in
the predicate's sense (pairwise transversal, every triple non-degenerate) has
its Clifford point at infinity, where the finite plane has no room for it. This is a statement about the
predicate, not a counterexample to Clifford's chain (Section~\ref{sec:chainfail}).
\item \textbf{Miquel's five-circles theorem} (Theorem~\ref{thm:five}).
\item \textbf{Lester's theorem} (Theorem~\ref{thm:lester}).
\item \textbf{The Sawayama--Th\'ebault theorem} (Theorem~\ref{thm:thebault}).
\item \textbf{Kosnita's theorem} (Theorem~\ref{thm:kosnita}).
\item \textbf{Blundon's inequalities} and \textbf{the Steiner--Lehmus theorem}
(Theorems~\ref{thm:blundon} and~\ref{thm:steiner}), which are not circle configurations but were
proved by the same method in the same project.
\end{enumerate}

\paragraph{Method.} The proofs follow the method of the earlier papers in this series
\cite{paperA,paperB}: choose coordinates (usually complex), reduce the geometric conclusion to the
vanishing of a polynomial in the coordinates, and prove that vanishing in Lean by an explicit
algebraic certificate (\lean{linear\_combination}, \lean{field\_simp}, \lean{nlinarith}); the exception
is the five-circles theorem, whose core lemma is an argument with doubled directed angles
(Section~\ref{sec:five}). The certificates come with
separate lemmas tying the coordinate objects to Mathlib's notions (circumcentre, incentre,
nine-point circle, \mdecl{Cospherical}, \mdecl{Concyclic}, \mdecl{Collinear}). They
are large (the $n=6$ Clifford case is the largest) but entirely mechanical; what needs care is the
\emph{statements}, and in particular their non-degeneracy hypotheses.

\paragraph{What is not claimed.} We do not prove Clifford's chain for general $n$, and we show that
the natural finite-plane statement of it is false for $n\ge6$ (Sections~\ref{sec:chain}
and~\ref{sec:chainfail}). That failure does not refute the textbook chain, in which every level is a
genuine circle or a unique point; the all-$n$ chain in that form remains open here
(Section~\ref{sec:l1415}). Every theorem is stated for configurations satisfying explicit
non-degeneracy hypotheses; Section~\ref{sec:degen} lists them and says which cannot simply be dropped.
Theorems are stated in the Euclidean plane, not in the inversive (M\"obius) plane, so ``circle'' means
a Euclidean circle, and lines appear only as stated alternatives.

\subsection{Availability}\label{sec:avail}
The proofs are in eight Lean~4 developments on Lean v4.33.1 with Mathlib pinned at revision
\texttt{0df444a3} (all eight lockfiles agree). File and line counts below are of the library
\lean{.lean} files, including the root import file, counting every physical line (comments and blank
lines included) and excluding scratch files; in total 24 files and 5{,}049 lines.
Supporting material (the verification battery reports, fresh-clone receipts, numerical scripts, blind re-formalisation files, reviews and hand-check sheet of version~1, and the Layer-15 reports, the blind writer's statements and report, the transcript audit and the reviews of version~2) is in the folder \texttt{PaperD-circle-configurations} of the repository \url{https://github.com/jpbald93/LEAN-verification}.
\begin{center}\small
\begin{tabular}{llrr}
\toprule
development & content & files & lines\\
\midrule
\lean{clifford} & Clifford $n=4,5,6$ & 3 & 1{,}271\\
\lean{cliffordchain} & the chain predicate: $n=4,5$; fails for $n\ge6$ & 7 & 1{,}679\\
\lean{fivecircles} & Miquel's five circles & 2 & 428\\
\lean{lester} & Lester's theorem & 4 & 506\\
\lean{thebault} & Sawayama--Th\'ebault & 2 & 559\\
\lean{kosnita} & Kosnita's theorem & 2 & 249\\
\lean{blundon}, \lean{steinerlehmus} & Blundon; Steiner--Lehmus & 2+2 & 228+129\\
\bottomrule
\end{tabular}
\end{center}
They are published on GitHub; the paper is checked against these commits, and every source reference
(\texttt{File}:\textit{line}) links to that line at the commit:
\begin{center}\small
\begin{tabular}{lll}
\toprule
development & repository & commit\\
\midrule
\lean{clifford} & \url{https://github.com/jpbald93/clifford} & \texttt{209abaa5e638}\\
\lean{cliffordchain} & \url{https://github.com/jpbald93/cliffordchain} & \texttt{ba083b2158e1}\\
\lean{fivecircles} & \url{https://github.com/jpbald93/fivecircles} & \texttt{f1fd9b7b2bd9}\\
\lean{lester} & \url{https://github.com/jpbald93/lester} & \texttt{9f94754987cb}\\
\lean{thebault} & \url{https://github.com/jpbald93/thebault} & \texttt{c37f92e53c9e}\\
\lean{kosnita} & \url{https://github.com/jpbald93/kosnita} & \texttt{702f0263a701}\\
\lean{blundon} & \url{https://github.com/jpbald93/blundon} & \texttt{118b4c9f3afd}\\
\lean{steinerlehmus} & \url{https://github.com/jpbald93/steinerlehmus} & \texttt{a3a3d26bee6c}\\
\bottomrule
\end{tabular}
\end{center}
Each development has a gate script that rejects \lean{sorry}, \lean{admit}, \lean{native\_decide},
\lean{axiom} and meta-programming commands in the sources, builds the library, and checks with
\lean{\#print axioms} that every listed declaration depends only on \lean{propext},
\lean{Classical.choice} and \lean{Quot.sound}; definitions and lemmas they depend on are covered
transitively. Every declaration cited with a qualified name in this paper, such as \decl{Clifford.clifford_one}, is in its development's list;
other names (in plain typewriter type) are auxiliary definitions or lemmas covered in that way.
The gate runs in its strict default mode (a permissive \lean{SCAFFOLD=1} mode exists for development
and was not used).

\emph{Version~2.} Each repository has one new published commit that adds two files, \lean{<Lib>/L14.lean} and
\lean{<Lib>/L15Classical.lean} (Section~\ref{sec:l1415}), two import lines in the root file, the
extended gate list and the README; the new commits are \lean{clifford} \texttt{927f7f82b3a1ae8342b1636b78bb9c04199c0d4b}, \lean{cliffordchain} \texttt{b6649f48bc4eaf192f729efd26a617fa2aeddfef}, \lean{fivecircles} \texttt{6040681ed19d180325f87162f690ebdad4a96961}, \lean{lester} \texttt{d307293ba23b92aece5f81b971704e6e11641bfa}, \lean{thebault} \texttt{2edce1eecc7b97faefa8580bddaef955ab44bf73}, \lean{kosnita} \texttt{7eaae5300ce68cf3c0eb59fbc52d8456c4c7f819}, \lean{blundon} \texttt{50e1f17b92bd09d2fc515fb46133f2f84b46ff6c} and \lean{steinerlehmus} \texttt{53b62463cb50388b51e7a703afeb08e7c8f1ebdc}. All other library files, and
the statements and source lines cited from them, are unchanged; the file and line counts above and
the source links are those of the commits in the table.

\section{Clifford's chain: $n = 4, 5, 6$}\label{sec:clifford}

\subsection{The model}
We identify the plane with $\CC$ and put the common point $M$ at the origin. A circle through $0$ is
determined by its centre $c\neq 0$: it is the circle $\lvert z-c\rvert=\lvert c\rvert$, i.e.\
$z\bar z = \bar c z + c\bar z$ (\decl{Clifford.circ_throughOrigin}). For two such circles with distinct centres
$c_i\neq c_j$, define
\[
 P_{ij} \;=\; \frac{\bar c_i\,c_j - c_i\,\bar c_j}{\bar c_i - \bar c_j}
\]
(\lean{secondInter}, \src{clifford}{Clifford}{Basic}{194}); for example $c_1=1$, $c_2=i$ give $P_{12}=1+i$.
Lean proves that $P_{ij}$ lies on both circles (\decl{Clifford.secondInter_mem_left}, \decl{Clifford.secondInter_mem_right}).
It is the second intersection point when it differs from $0$; it equals $0$ exactly when the circles
are tangent at $0$ (e.g.\ $c_1=1$, $c_2=2$), which the theorems below exclude by assuming $P_{ij}\neq 0$.
For three circles $i,j,k$, $\Gamma_{ijk}$ (\lean{gam}, \src{clifford}{Clifford}{Basic}{203}) is the
circumcircle of $P_{ij},P_{ik},P_{jk}$, with centre given by an explicit circumcentre formula
(\lean{ccenter}, \src{clifford}{Clifford}{Basic}{197}); when the three points are not
collinear Lean proves that they lie on it (\decl{Clifford.mem_circumcircle_left} and its two
companions). Mathlib's \mdecl{EuclideanGeometry.Sphere} is used for circles throughout.

\subsection{The theorems}

\begin{theorem}[Clifford, $n=4$; \decl{Clifford.clifford_one}]\label{thm:cliff4}
Let four circles pass through a common point $M$, with centres $c_1,\dots,c_4$. Suppose every
second intersection point $P_{ij}$ ($i<j$) differs from $M$, and for each triple $\{i,j,k\}$ the
three points $P_{ij},P_{ik},P_{jk}$ are not collinear. Then the four circles
$\Gamma_{123},\Gamma_{124},\Gamma_{134},\Gamma_{234}$ have a common point.
\end{theorem}

The theorem (\src{clifford}{Clifford}{Basic}{381}) asserts existence only. Its witness is an explicit
point, a ratio of two determinant-like polynomials in the centres and their conjugates,
$\mathrm{cp}(c_1,c_2,c_3,c_4)$ (\lean{cliffordPoint}, \src{clifford}{Clifford}{Basic}{309}); Lean proves
that it lies on each $\Gamma_{ijk}$ (\decl{Clifford.cliffordPoint_mem_gam}, \src{clifford}{Clifford}{Basic}{323},
and \decl{Clifford.cliffordPoint_mem_gams} for all four at once). We do not prove that the common point
is unique, and the theorems below are about these explicit points.

\begin{theorem}[Clifford, $n=5$; \decl{Clifford.clifford_two}, \decl{Clifford.clifford_two_explicit}]\label{thm:cliff5}
Let five circles pass through $M$, with every $P_{ij}\neq M$ and no triple $P_{ij},P_{ik},P_{jk}$
collinear. For each $k$, let $P_k=\mathrm{cp}$ of the four circles other than circle $k$. Then each
$P_k$ lies on the four circles $\Gamma$ of those four, and the five points $P_1,\dots,P_5$ lie on one
circle or on one line.
\end{theorem}

In Lean the statement is split in two. \decl{Clifford.clifford_two} (\src{clifford}{Clifford}{Third}{543})
asserts that there \emph{exist} five points with these incidences that lie on a circle or a line;
\decl{Clifford.clifford_two_explicit} (\src{clifford}{Clifford}{Third}{457}) proves the circle-or-line
conclusion for the five explicit points $\mathrm{cp}$, and \decl{Clifford.cliffordPoint_mem_gams} gives
their incidences. The disjunction is not a weakening: in the independent agent's numerical check (reported figures, script not kept), configurations
satisfying every hypothesis were found in which the five points are pairwise distinct (minimum spacing
$0.042$) and collinear to $5\times10^{-60}$. The second alternative corresponds to the usual inversive
formulation, in which a line is a circle through $\infty$; our statement is the finite Euclidean
circle-or-line form of it, under the stated hypotheses.

\begin{theorem}[Clifford, $n=6$; \decl{Clifford.clifford_three}]\label{thm:cliff6}
Let six circles pass through $M$, with every $P_{ij}\neq M$ and no triple $P_{ij},P_{ik},P_{jk}$
collinear. For each five-element subset $S$ of the six, let $P^S_T=\mathrm{cp}(c_T)$ for the five
four-element subsets $T\subset S$, and let $C_S$ be the circle through three designated ones of them
(for $S=\{1,\dots,5\}$: those of $\{1,2,3,4\}$, $\{1,2,3,5\}$, $\{1,2,4,5\}$; the six choices are listed
in the Lean statement). Suppose those three points are not collinear, for each $S$. Then all five
points $P^S_T$ lie on $C_S$, for each $S$, and the six circles $C_S$ have a common point.
\end{theorem}

The six extra hypotheses (\lean{m1}--\lean{m6} in \src{clifford}{Clifford}{Third}{780}) say that the
three designated points of each $S$ are not collinear, so that $C_S$ is a genuine circle; they exclude
exactly the case that Theorem~\ref{thm:cliff5} allows in its second alternative (and the case of
coincident points). So Theorem~\ref{thm:cliff6} is narrower than the
generalized-circle statement: it is Clifford's third theorem for configurations whose five-point
circles are honest circles. The proof reduces all six memberships to one polynomial identity
(\decl{Clifford.master}, \src{clifford}{Clifford}{Third}{187}), the largest certificate in this paper.

\subsection{The chain for general $n$}\label{sec:chain}
A second, independently written development, \lean{cliffordchain}, defines an existential incidence
predicate for the whole chain. Its \emph{generalized circles} (\lean{IsGCircle},
\src{cliffordchain}{Clifford}{Clean}{156}) are lines and Mathlib spheres \mdecl{Metric.sphere}~$o\,r$ for
any real $r$; so a radius-$0$ ``circle'' (a single point) and, formally, the empty set ($r<0$) are
included. The proved case $n=5$ uses the radius-$0$ case when all five points coincide. For centres $c:\{0,\dots,n-1\}\to\CC$, the predicate
$\mathrm{Chain}(n,c,X)$ (\src{cliffordchain}{Clifford}{Clean}{165}) says that $X$ is a Clifford object of
level $n$: for $n=0,1,2$ the point $M$, the circle itself, and the point $P_{01}$; for odd $n\ge3$, a
generalized circle containing, for each $i$, \emph{some} level-$(n-1)$ point of the circles other than
$i$; for even $n\ge4$, a single point lying, for each $i$, on \emph{some} level-$(n-1)$ object of the
circles other than $i$. The predicate asserts existence at each level; it does not assert that the
objects are unique or that the choices at different levels cohere. The development's intended general statement is $\mathrm{CliffordChain}(n)$: in general
position (pairwise transversal, every triple non-degenerate) a level-$n$ object exists
(\src{cliffordchain}{Clifford}{Clean}{184}). The general-position conditions are stated as
polynomial inequalities, and a bridge file ties them to geometry (\src{cliffordchain}{Clifford}{Bridge}{58}
onwards): \lean{Transversal}~$c\,d$ holds if and only if $M$, $c$, $d$ are not collinear
(\decl{Clifford.transversal_iff}); the first condition of \lean{GenTriple} holds if and only if the three
centres are not collinear (\decl{Clifford.triP_ne_zero_iff}); and, for pairwise transversal circles with
non-collinear centres, the second holds if and only if the three second intersection points are not
collinear (\decl{Clifford.cycP_ne_zero_iff}). The last needs the non-collinear centres: the underlying
identity is $X\prod\lvert c_i-c_j\rvert^2=\mathrm{triP}^2\cdot\mathrm{cycP}$, where $X$ is the collinearity
determinant of the three points. Together, \decl{Clifford.genTriple_iff} says that \lean{GenTriple} means
exactly ``neither the centres nor the three second intersection points are collinear''; and
\decl{Clifford.onCircleM_iff_mem_sphere} identifies the circle model with Mathlib's \mdecl{Metric.sphere}.

\begin{theorem}[\decl{Clifford.cliffordChain_four}, \decl{Clifford.cliffordChain_five}]\label{thm:chain}
$\mathrm{CliffordChain}(4)$ and $\mathrm{CliffordChain}(5)$ hold.
\end{theorem}

The $n=5$ case rests on \decl{Clifford.clifford_two_clean} (\src{cliffordchain}{Clifford}{Clean}{123}):
in general position the five points are concyclic or collinear, with no further hypothesis. (The
same development also contains a version of the $n=5$ theorem with an extra coefficient condition,
\lean{CliffordNondeg}, which excludes some genuine configurations; we do not cite that version as
Clifford's theorem.) For $n\ge6$ the predicate is in fact false; see Section~\ref{sec:chainfail}.

\subsection{The finite-plane chain fails for $n\ge6$}\label{sec:chainfail}
The predicate $\mathrm{Chain}$ places every Clifford object in the finite plane $\CC$. Classically,
Clifford's theorem lives in the inversive plane: after inversion at $M$ the $n$ circles become $n$ lines
in general position, whose Clifford objects are honest points and circles, and the inverse image of a
Clifford point that happens to be $M$ itself is the point at infinity. Nothing in ``pairwise transversal,
every triple non-degenerate'' prevents this, and the finite-plane predicate has no point at infinity to
use. We prove that this already happens for six circles, and therefore for every $n\ge6$.

\begin{theorem}[\decl{Clifford.Counter6.not_cliffordChain_six}, \decl{Clifford.not_cliffordChain_of_six_le}, \decl{Clifford.cliffordChain_iff}]\label{thm:chainfail}
For the six centres
\[
  c=(3-6i,\;6-3i,\;1+2i,\;-2-i,\;-4i,\;4)
\]
the circles through $M=0$ are in general position in the sense of $\mathrm{CliffordChain}$, but no
level-$6$ object exists. Consequently $\mathrm{CliffordChain}(n)$ fails for every $n\ge6$, and
$\mathrm{CliffordChain}(n)$ holds if and only if $n\le5$.
\end{theorem}

Number the centres $c_1,\dots,c_6$ in the order displayed. For this configuration every level-$4$
object is forced (by the three circles through it) to be the classical Clifford point, and for each
of the six five-circle subfamilies the five such points are distinct and \emph{collinear}, so the
level-$5$ objects are forced to be lines: for instance $\operatorname{Re}z=4$ when circle~$5$ is omitted
and $\operatorname{Im}z=-4$ when circle~$6$ is omitted. These two lines meet at $4-4i$, which is not on
the line for omitted circle~$1$, so no point lies on all three, let alone all six, level-$5$ objects
(\src{cliffordchain}{Clifford}{Counter6}{182}). The kernel checks what the proof needs: general position,
the nine level-$4$ points that determine these three lines, and their non-concurrence, all in exact
arithmetic over $\mathbb{Q}(i)$. The full picture (all fifteen level-$4$ points and all six lines) was
checked outside Lean in exact rational arithmetic, by the agent and independently by us. Replacing $c_6=4$ by $4+\varepsilon$ (real $\varepsilon\ne0$, small enough that the configuration stays in general position) gives, in exact arithmetic, the
classical level-$6$ point $4-4i+50/\varepsilon$, which goes to infinity as $\varepsilon\to0$. The step to $n\ge7$ uses that a level-$(n+1)$ object requires a level-$n$
object of the subfamily omitting the last circle, together with the fact that any family in general
position can be extended by one more circle in general position (the new centre is taken on a parabola,
which meets each excluded locus in finitely many points); see \src{cliffordchain}{Clifford}{Classify}{57}
and \src{cliffordchain}{Clifford}{Classify}{167}.

Together, Theorems~\ref{thm:chain} and~\ref{thm:chainfail} settle the predicate completely: it is a
correct formalization of the cases $n=4,5$ but not of Clifford's theorem for $n\ge6$. The failure is
about this predicate, in which odd levels may be a circle \emph{or a line} and general position is
imposed only on the initial pairs and triples. It does not refute the textbook chain, in which every
level is a genuine circle, or a point that is the unique common point of its circles: Layer~15 proves
(\decl{Clifford.Classical.cov_D1_d_no_spherical_level_five}) that for the six centres above any level-$5$
object of the subfamily omitting circle~$6$ contains the three distinct collinear points $2/3-4i$,
$-8/3-4i$, $-6-4i$ and so lies on no circle; hence a set whose complex image is such an object lies on no
circle of the textbook formulation (\decl{Clifford.Classical.cov_D1_d_counterexample_not_writer_data}, which
takes this identification as a premise; the full translation between the two chain structures is not
formalized). So this configuration cannot supply the genuine level-$5$ circles that the textbook hypotheses require. A faithful general statement needs either
the Riemann sphere (even-level objects in $\CC\cup\{\infty\}$, lines counted as circles through
$\infty$) or the equivalent formulation for $n$ lines. In the same exploration (60 digits, random
centres) the classical chain exists for $n\le8$ and agrees with a closed form after inversion, but we
have not formalized either corrected statement. The six-circle theorem of Section~\ref{sec:clifford}
(Theorem~\ref{thm:cliff6}) is unaffected: its hypotheses $m_1$--$m_6$ require the designated triples to be non-collinear, and in this
configuration every five-point set is collinear (an argument, not a machine-checked statement; for the
subfamily omitting circle~$6$, Layer~15 proves that any level-$5$ object contains three distinct collinear
points, Section~\ref{sec:l1415}).

\section{Miquel's five circles}\label{sec:five}

Take a pentagon $A_0,\dots,A_4$, extend its sides to form a pentagram with outer points
$T_0,\dots,T_4$, and draw the five circles through $A_{i-1}$, $A_i$ and the triangle point $T_i$.
Consecutive circles meet again at points $M_i$. Miquel's theorem \cite{miquel1838} says that the five
points $M_i$ are concyclic.

\begin{theorem}[\decl{Fivecircles.fivecircles}]\label{thm:five}
Let $A,T,M:\{0,\dots,4\}\to\CC$ be injective, with no point of one family equal to a point of
another (so the fifteen labelled points are pairwise distinct). Suppose that for each $i$ (indices mod~5) the points $A_{i-1},A_i,T_{i+1},T_{i-1}$ are
collinear and the points $A_{i-1},A_i,T_i,M_{i-1},M_i$ are concyclic. Then $M_0,\dots,M_4$ lie on one
circle or on one line.
\end{theorem}

The collinearity hypotheses say that $T_{i\pm1}$ lie on the line $A_{i-1}A_i$, which is the
pentagram incidence; the concyclicity hypotheses say that $M_{i-1}$ and $M_i$ are the second
intersections of the circle of triangle $i$ with its neighbours. Distinctness of the fifteen points
excludes, in particular, tangent circles (a second intersection coinciding with $A_i$). Again the
alternative is genuine: the independent agent's numerical check (reported figures, script not kept) found a non-degenerate pentagram (minimum point
separation $0.52$) whose five points $M_i$ are collinear to $10^{-39}$. Unlike the other developments,
the core of this proof is not polynomial algebra. Concyclicity of four points is reduced to the
vanishing of the imaginary part of a cross-ratio, and that to an equality of doubled directed angles
(\mdecl{Real.Angle}); the core lemma (\decl{Fivecircles.core}, \src{fivecircles}{Fivecircles}{Basic}{167})
is then an angle chase, written as a sequence of additive certificates, and is reused for the five
cyclic relabellings. The theorem is \src{fivecircles}{Fivecircles}{Basic}{383}.

\section{Lester's circle}\label{sec:lester}

\begin{theorem}[Lester; \decl{Lester.lester}]\label{thm:lester}
Let $ABC$ be a scalene triangle in the plane. Let $F_1$ be its first Fermat point, the common point of
the three lines joining each vertex to the apex of the equilateral triangle erected outwards on the
opposite side, and $F_2$ its second Fermat point (the same with the equilateral triangles erected
inwards). Then the circumcentre, the nine-point centre, $F_1$ and $F_2$ are concyclic.
\end{theorem}

The theorem is due to Lester \cite{lester1997}; see \cite{yiu2010} for the circles of Lester, Evans
and Parry. The hypothesis is \emph{scalene}, not merely ``not equilateral''. For an isosceles triangle
all four points lie on the axis of symmetry, and for a typical one they are four distinct points: for
$A=(0,h)$, $B=(-1,0)$, $C=(1,0)$ with $h\in\{0.3,\,0.5,\,1,\,2.5\}$ we computed (60 digits, script in the verification repository) four distinct collinear
points (minimum separation $0.048$ to $0.33$), which are not concyclic. So the scalene hypothesis cannot
simply be weakened to ``not equilateral''. It does more than is needed, though: in some isosceles
triangles points coincide and the conclusion holds trivially, since \mdecl{Concyclic} is a property of
a \emph{set}. For the isosceles triangle with apex angle $120^\circ$, $A=(0,1/\sqrt3)$, the nine-point
centre and the first Fermat point both equal $A$, and the second Fermat point equals the circumcentre,
so the set of four points has only two elements.
The circumcentre is Mathlib's; the nine-point centre is defined as the midpoint of the circumcentre and
the orthocentre and proved equal to the centre of Mathlib's nine-point circle
(\decl{Lester.bridge_ninePointCenter}); the equilateral apexes are proved equilateral
(\decl{Lester.bridge_apexOut}, \decl{Lester.bridge_apexIn}); and $e^{i\pi/3}$ is identified with the
Cartesian constant used (\decl{Lester.bridge_ω}). The theorem is \src{lester}{Lester}{Basic}{56}. It
assumes the Fermat points and does not assert their existence; existence, for scalene triangles and
with the Fermat points defined geometrically, is proved in the Layer-15 file
(\decl{Lester.Classical.cov_D3_obligation}, Section~\ref{sec:l1415}).

\section{The Sawayama--Th\'ebault theorem}\label{sec:thebault}

Th\'ebault proposed in 1938 \cite{thebault1938} the following problem, which had been proved already
by Sawayama in 1905 \cite{sawayama1905}; see \cite{ayme2003} for its history.

\begin{theorem}[Sawayama--Th\'ebault; \decl{Thebault3.sawayama_thebault}]\label{thm:thebault}
Let $ABC$ be a triangle and $D$ a point strictly between $B$ and $C$. Let $P$ be the centre of a
circle tangent to the lines $AD$ and $BC$, internally tangent to the circumcircle of $ABC$, and lying in
the angle $\angle ADB$; let $Q$ be the centre of such a circle lying in the angle $\angle ADC$. Then
$P$, $Q$ and the incentre of $ABC$ are collinear.
\end{theorem}

The definition \lean{IsThebaultCircle} (\src{thebault}{Thebault3}{Basic}{32}) says that there is
$r>0$ with $P$ at distance $r$ from both lines, at distance $R-r$ from the circumcentre, on the same
side of $AD$ as $B$ (resp.\ $C$) and on the same side of $BC$ as $A$; the two side conditions
together are exactly ``in the open angle''. In the independent numerical check, each angle
contained exactly one such circle, so the definition picks out the intended circle. The incentre is Mathlib's, and
\decl{Thebault3.tri_incenter_spec} proves it equidistant from the three side lines and on the inner
side of each. ``Exactly one such circle'' is the agent's reported observation on 1989 sampled triangles (script not kept), not a proved
uniqueness statement. The theorem assumes $P$ and $Q$ and does not assert their existence; the
Layer-15 file proves existence of two such circles in its own formulation (tangent to both lines,
internally tangent to the circumcircle, centres on opposite sides of $AD$ and on $A$'s side of $BC$;
\decl{Thebault3.Classical.cov_D4_obligation}, Section~\ref{sec:l1415}). It is
\src{thebault}{Thebault3}{Basic}{514}.

\section{Kosnita's theorem}\label{sec:kosnita}

\begin{theorem}[Kosnita; \decl{Kosnita.kosnita}]\label{thm:kosnita}
Let $ABC$ be a triangle with circumcentre $O$, and let $O_a,O_b,O_c$ be the circumcentres of the
triangles $OBC$, $OCA$, $OAB$. Suppose that no two of $A,B,C$ are diametrically opposite (so that
$OBC$, $OCA$, $OAB$ are triangles) and that $AB^2+BC^2+CA^2\neq 5R^2$, where $R$ is the circumradius.
Then the lines $AO_a$, $BO_b$, $CO_c$ are concurrent.
\end{theorem}

The circumcentres are specified by equidistance (\lean{hOa1}, \lean{hOa2}, \ldots), which in the
plane is equivalent. The first exclusion (\lean{hOBC}, \lean{hOCA}, \lean{hOAB}: $O,B,C$ not collinear, etc.) excludes
right triangles, for which $O$ is the midpoint of a side and $O_a$ does not exist. In the Lean
statement it is in fact redundant: if $O,B,C$ were collinear, then $O$ would be the midpoint of $BC$, and
no point is equidistant from three distinct collinear points, so the equidistance hypotheses for $O_a$
could not hold. We keep it because the proof uses it. The second, $\lean{hS}$, cannot simply be
dropped. The concurrence point (classically the Kosnita point, $X_{54}$ in Kimberling's encyclopedia
\cite{etc}; our theorem does not identify it) can fail to exist on the locus $AB^2+BC^2+CA^2=5R^2$.
In the independent agent's numerical check, three sampled triangles on that locus
gave three pairwise parallel, distinct lines (determinants of the line pairs about $10^{-61}$,
separations about $0.18$; reported figures, script not kept), so for them no common point exists and
the conclusion fails. Not every triangle on the locus behaves this way. For the isosceles triangle $A=(1,0)$, $B,C=(\tfrac12,\pm\tfrac{\sqrt3}2)$ on the unit circle
(sides $1,\sqrt3,1$, angle $120^\circ$ at $A$), $O_a=A$, so ``the line $AO_a$'' is undefined, while $BO_b$ and $CO_c$ are the same line; here the Lean conclusion holds trivially (any point
of that line will do). So we do not claim a uniform description of the excluded locus; only that the
affine conclusion fails on part of it. Relative to a projective statement ``concurrent, possibly at infinity'' the exclusion is a
weakening. Layer~15 removes it (Section~\ref{sec:l1415}): from the same polynomial certificate it proves
that on the locus $AB^2+BC^2+CA^2=5R^2$ the three lines, when defined ($A\neq O_a$ etc.), are pairwise
parallel (\decl{Kosnita.Classical.cov_Kosnita_zero_locus}), so ``concurrent or all parallel'' holds for
every triangle for which the three lines are defined (\decl{Kosnita.Classical.cov_D5}). The
theorem is \src{kosnita}{Kosnita}{Basic}{180}; see \cite{grinberg2003} for the Kosnita point.

\section{Two companions: Blundon and Steiner--Lehmus}\label{sec:companions}

\begin{theorem}[Blundon; \decl{Blundon.blundon}]\label{thm:blundon}
For a triangle in a real inner product space with circumradius $R$, inradius $r$ and semiperimeter $s$,
\[
2R^2+10Rr-r^2-2(R-2r)\sqrt{R^2-2Rr}\;\le\; s^2\;\le\;2R^2+10Rr-r^2+2(R-2r)\sqrt{R^2-2Rr}.
\]
\end{theorem}
Blundon proved that these bounds are best possible \cite{blundon1965}. $R$ and $r$ are Mathlib's
\mdecl{Affine.Simplex.circumradius} and \mdecl{Affine.Simplex.inradius}. The proof is a single
polynomial certificate (\decl{Blundon.core_identity}): with $x,y,z$ the tangent lengths,
$s^3\bigl(4r^2E+\Delta^2\bigr)=0$ where $E=s^4-2(2R^2+10Rr-r^2)s^2+r(4R+r)^3$ and $\Delta$ is the product
of the differences of the sides; hence $E\le 0$, which is the stated pair of inequalities. Bridges tie
$R$ and $r$ to the side lengths (\decl{Blundon.circumradius_sq_mul}, \decl{Blundon.height_sq_mul}).
The theorem is \src{blundon}{Blundon}{Basic}{147}.

\begin{theorem}[Steiner--Lehmus; \decl{SteinerLehmus.steiner_lehmus}]\label{thm:steiner}
In a non-degenerate triangle $ABC$, let $E$ on segment $AC$ and $F$ on segment $AB$ be the feet of
the internal bisectors from $B$ and $C$. If $BE=CF$ then $AB=AC$.
\end{theorem}
The bisector-length formula $BE^2(BA+BC)^2 = BA\cdot BC\,((BA+BC)^2-AC^2)$ is proved from Mathlib's
angles (\decl{SteinerLehmus.bisector_foot_sq}); the two cleared formulas then force $b=c$ through the
factorisation $(b-c)\cdot a(a+b+c)(a^3+a^2b+a^2c+3abc+b^2c+bc^2)=0$
(\decl{SteinerLehmus.algebraic_core}). The theorem holds in any real inner product space, and is
\src{steinerlehmus}{SteinerLehmus}{Basic}{112}. (The proof is direct; the theorem's notoriety for
lacking a direct proof \cite{gilbert1963,coxeter1967} concerns synthetic proofs.) The algebraic step
alone, $b=c$ from equality of the two bisector-length \emph{formulas} over positive side lengths, has
an earlier sorry-free Lean proof \cite{leangeniusSL}; what is new here is the geometric statement, with
the bisector feet given as points and the bisector condition as an equality of Mathlib angles, and the
bridge \decl{SteinerLehmus.bisector_foot_sq} from it to the formula.

\section{Non-degeneracy: what each statement excludes}\label{sec:degen}

\begin{center}\small
\begin{tabular}{p{2.6cm}p{6.2cm}p{5.2cm}}
\toprule
theorem & excluded configurations & status\\
\midrule
Clifford $n=4,5,6$ & tangent circles or equal centres ($P_{ij}=M$); a $\Gamma_{ijk}$ that is a line
(three collinear $P$'s) & standard general position\\
Clifford $n=5$ & --- (conclusion: circle \emph{or} line) & line case genuine\\
Clifford $n=6$ & three designated points of some $C_S$ collinear or coincident & narrower than the inversive statement\\
Five circles & any coincidence among the fifteen labelled points, incl.\ tangent circles & line case genuine\\
Lester & isosceles triangles; Fermat points assumed (existence: Layer~15) & cannot be weakened to ``not equilateral''\\
Sawayama--Th\'ebault & $D$ at $B$ or $C$; degenerate triangle & classical setting\\
Kosnita & right triangles (redundant given the other hypotheses); $AB^2+BC^2+CA^2=5R^2$ & $5R^2$ case cannot be dropped (affine); lines parallel there (Layer~15)\\
Blundon, Steiner--Lehmus & degenerate triangle & ---\\
\bottomrule
\end{tabular}
\end{center}

\section{Does the statement mean what it says?}\label{sec:blind}
The verification battery checks that proofs are accepted and gives partial evidence against vacuous
statements; it does not
show that a Lean statement says what the prose says. To test that, we removed every code name, code
listing and link from the mathematical sections of the paper and gave the remaining prose to a
language model that had not worked on the developments, with instructions to read nothing else. (Its
tool log, in the verification repository, shows that the only file it read was the prepared text; it also compiled a short
probe of Mathlib names, and its own statements, on the x86-64 machine; the text was produced before
Section~\ref{sec:chainfail} was written, so Theorem~\ref{thm:chainfail} is not covered.) It wrote twelve
Lean statements: one for each theorem, two for Theorem~\ref{thm:cliff5} (the existential and the explicit
form) and two for Theorem~\ref{thm:chain} ($n=4,5$), with every notion it needed defined from the prose,
and it marked the places where it found the prose ambiguous. A second model, different from the
writer, then tried to prove in Lean, for each pair, that the published statement implies the blind
one and conversely. In these proofs it could use the premise, unfolding of definitions, and general
Mathlib lemmas, together with the small bridge lemmas about the developments' own notions (for example
\decl{Clifford.transversal_iff} and \decl{Lester.bridge_apexOut}), but never the published theorem
itself, so that the check is about the statements and not about their truth. We recompiled the
comparison files ourselves (no \lean{sorry}, no axioms) and checked that no published theorem is used
in a proof.

Ten of the twelve pairs are equivalent in both directions: Theorems~\ref{thm:cliff4} and~\ref{thm:chain}
(both parts), the existential part of Theorem~\ref{thm:cliff5}, and Theorems~\ref{thm:five}
to~\ref{thm:steiner}. The equivalences are not all literal. Lester's theorem was stated over
$\mathbb{R}^2$ with apexes defined geometrically and is proved in $\mathbb{C}$ with explicit apexes; the
blind Kosnita statement uses lines and ``$O$ is not the midpoint of a side'' where ours uses collinear
triples and non-collinearity; the blind chain predicate is defined by recursion on objects rather than
sets. The two remaining pairs are the explicit-point parts of Theorems~\ref{thm:cliff5}
and~\ref{thm:cliff6}. Our prose does not give the formula of the explicit point \lean{cp}, so the blind
writer quantified over \emph{every} common point of the four circles, as its ambiguity note records.
Each direction then holds once one missing fact is supplied as a premise: that \lean{cp} is a common
point (from blind to ours), or that the common point is unique (from ours to blind). We have not proved
uniqueness. So for these two theorems the check found a real gap, though in our prose rather than in
the Lean: the prose does not determine the Lean statement. Theorems~\ref{thm:cliff5} and~\ref{thm:cliff6}
are about the explicit point $\mathrm{cp}$ of Section~\ref{sec:clifford}, whose defining formula is in
the source and not in the text; read as statements about every common point, they would claim more than
we prove. For Theorem~\ref{thm:cliff6} the writer also had to guess which three points of each
five-point set are designated; its guess agrees with ours for all six sets. Apart from this
underspecification, no difference between prose and Lean was found. The blind text, the blind
statements, the writer's tool log, the comparison files, our recompilation log and both models' reports
are in the verification repository.

\section{Related work}\label{sec:related}

\paragraph{Formal proofs.} The Archive of Formal Proofs entry on Miquel's theorem \cite{afpmiquel} proves both the incidence
form of the three-circle (pivot) theorem (a point on two of the circles lies on the third) and the
existence of the common point; Rocq's HighSchoolGeometry \cite{hsgeometry} proves the incidence form of
the four-line Miquel theorem; and a community Lean repository \cite{leangenius} proves the incidence
form of the pivot theorem. These three served as a control for our searches.
Mathlib's list of 1000 theorems \cite{mathlib1000} lists Clifford's circle theorems, the five circles
theorem, Lester's theorem, Th\'ebault's theorem and the Steiner--Lehmus theorem, all without a
formalisation; Kosnita and Blundon are not on the list.

On 13 October 2026 we searched for each of the classical theorems formalised here
(Theorems~\ref{thm:cliff4}--\ref{thm:chain} and~\ref{thm:five}--\ref{thm:steiner}; the predicate-specific
Theorem~\ref{thm:chainfail} was not a search target): Mathlib
(code, \texttt{docs/100.yaml} and \texttt{docs/1000.yaml}, issues and pull requests); the full
Archive of Formal Proofs \cite{afp} (a mirror of 7 October 2026 with 1045 entries, by full-text search,
plus the live front page for newer entries);
all 1500 articles of the Mizar library; full clones of HOL4 and HOL Light; Metamath's \texttt{set.mm};
Rocq's GeoCoq and HighSchoolGeometry; ProofAtlas (sitemap and search index); the public archive of
the Lean Zulip; compfiles and \texttt{formal-conjectures}; the lean-genius repository \cite{leangenius}
(full tree and pull requests); and GitHub code search. In this scope we found no formalisation,
complete or partial,
of Theorems~\ref{thm:cliff4}--\ref{thm:chain} and~\ref{thm:five}--\ref{thm:kosnita}, and no complete formal proof of Blundon's
inequalities; the nearest is a Lean proof of Gerretsen's weaker inequality $s^2\le4R^2+4Rr+3r^2$
\cite{leangeniusG}. For Steiner--Lehmus the algebraic core is proved in Lean \cite{leangeniusSL}
(Section~\ref{sec:companions}), and the Acorn library states the theorem without a proof
\cite{acornSL}.

The search has limits. The Zulip search covered the public archive only, up to about 2
October 2026; the live search needs a login. GitHub's code index does not cover every repository,
returns at most about 50 results per query, and rate limits forced us to use queries not restricted by
language, which we filtered by hand. ProofAtlas has no public API, so we used its own index and a
third-party \lean{site:} search. The absence of a hit is a statement about this scope only.

\paragraph{Machine proofs outside proof assistants.} Wu's method and its successors prove many of
geometry theorems by computer algebra, among them Th\'ebault's problem \cite{chou1988}. Java Geometry
Expert lists Steiner--Lehmus among its examples, and a Gr\"obner-basis prover (Hasclid) has an example
file for a Blundon-type bound, which we did not run. None of these is checked by a proof kernel. In
Wu's method the non-degeneracy conditions are produced by the method \cite{chou1988}; our statements
state them explicitly.

\paragraph{Our earlier papers.} Papers~A \cite{paperA} and~B \cite{paperB} formalised other classical
results by the same coordinate-certificate method, and Paper~C \cite{paperC} treated Barbier's theorem
and the four-vertex theorem.

\section{Verification}\label{sec:verify}

\paragraph{Gates and tampering tests.} The gates pass with only the three standard axioms:
\lean{clifford} 20 declarations, \lean{cliffordchain} 32, \lean{lester} 19, \lean{thebault} 12,
\lean{kosnita} 11, \lean{fivecircles} 6, \lean{blundon} 5, \lean{steinerlehmus} 3 (108 in all). Each
tampering suite plants a hidden \lean{sorry}, an \lean{axiom} split across two lines, a fake
\lean{\#print axioms} report, a command-redefining macro, a wrapper axiom, and one or more broken
hypotheses or bridges that should make the build fail; every planted defect was rejected.
The table records the final receipts. On 14 October 2026, on each machine, a script cloned each of the
eight repositories from GitHub, checked out the commit pinned in this paper (and recorded that the
working tree was clean), recorded the Lean and Mathlib revisions and the SHA-256 hash of every library
file, built the library with only the pinned Mathlib supplied, and ran the gate and the tampering suite.
The two logs are in the verification repository.
\begin{center}\small
\begin{tabular}{lcc}
\toprule
development & x86-64 Linux & arm64 macOS\\
\midrule
\lean{clifford} & 20 / tamper OK & 20 / tamper OK\\
\lean{cliffordchain} & 32 / tamper OK & 32 / tamper OK\\
\lean{fivecircles} & 6 / tamper OK & 6 / tamper OK\\
\lean{lester} & 19 / tamper OK & 19 / tamper OK\\
\lean{thebault} & 12 / tamper OK & 12 / tamper OK\\
\lean{kosnita} & 11 / tamper OK & 11 / tamper OK\\
\lean{blundon} & 5 / tamper OK & 5 / tamper OK\\
\lean{steinerlehmus} & 3 / tamper OK & 3 / tamper OK\\
\bottomrule
\end{tabular}
\end{center}
Entries are the number of declarations for which the gate passed, and whether every planted defect was
rejected. All sixteen checkouts were at the pinned commit with a clean working tree; the 24 library-file hashes
agree between the two logs and with the published commits, and every build used Mathlib
\texttt{0df444a3} (recorded separately, because a bug in the script left that field of the logs
empty). In the first run the script checked out the previous commit of \lean{cliffordchain}, which
differs from the pinned one only in a text file; that development was re-run at the pinned commit on
both machines with a corrected script, with the results shown, and both runs are in the verification repository. (Earlier gate runs, on intermediate versions, are also there but are superseded by these.)

\paragraph{Independent re-verification.} For six of the eight developments (both Clifford
developments, \lean{fivecircles}, \lean{lester}, \lean{thebault}, \lean{kosnita}) a separate agent, not
the one that wrote the proofs, read the statements against the classical theorems hypothesis by
hypothesis and re-checked them numerically; its reports are in the verification repository. Their scripts were not kept, so
we re-did the numerical part reproducibly for all eight developments, with a single script that computes each configuration from coordinate formulas (for Lester it
follows the Lean development's orientation and apex conventions, so it cannot detect a convention
error there; the conclusion is symmetric in $F_1,F_2$ anyway), fixed random seeds, 60-digit
arithmetic, and, for every theorem except Steiner--Lehmus, a deliberately false variant as a negative
control on the same configurations (script,
seeds, output and SHA-256 checksums in the verification repository). Results, as configurations / worst residual /
false variant violated:
\begin{center}\small
\begin{tabular}{lrrl}
\toprule
theorem & configs & worst residual & negative control violated\\
\midrule
Clifford $n=4$ & 1000 & $1.1\times10^{-58}$ & 1000/1000\\
Clifford $n=5$ & 500 & $1.2\times10^{-57}$ & 500/500\\
Clifford $n=6$ & 100 & $1.1\times10^{-57}$ & 100/100\\
Five circles & 500 & $1.8\times10^{-59}$ & 500/500\\
Lester & 1000 & $1.6\times10^{-58}$ & 1000/1000\\
Sawayama--Th\'ebault & 300 & $2.4\times10^{-60}$ & 300/300\\
Kosnita & 1000 & $1.2\times10^{-59}$ & 1000/1000\\
Steiner--Lehmus & 200 & $7.8\times10^{-62}$ & ---\\
Blundon & 1000 & (min.\ margin $4.5\times10^{-11}\ge0$) & 310/1000\\
\bottomrule
\end{tabular}
\end{center}
Each configuration satisfies the hypotheses of the Lean statement up to a margin (pairwise and triple
non-degeneracy of the input circles, and of the output triples the script tests, at least $10^{-3}$
for Clifford; the designated triples of Theorem~\ref{thm:cliff6} are non-collinear, being distinct concyclic
points, but are not held to that margin; scalene by at least $10^{-2}$ for Lester, and so on). The
Steiner--Lehmus configurations were found by solving $BE=CF$ for the third vertex from random starting
points; every solution found was isosceles. For Blundon the residual column gives the smallest relative
margin by which $s^2$ lies inside the two bounds, and the control is the false bound obtained by dropping
the square-root term, which fails in 310 of 1000 triangles. (In the agents' earlier runs, one five-circles control held in 1 of 1000 trials; in the re-run
every control except Blundon's was violated in every trial. For Clifford $n=6$ the script checks
non-collinearity of three of the five points of each $C_S$, not of the three designated in the Lean
statement; the designated ones follow, since the five points are distinct and concyclic.) Numerical
checks are evidence, not proofs.

\paragraph{Statement-level checks.} Every development was then run through the statement-level
battery used in the earlier papers, on 13 and 14 October 2026: \lean{clifford} on the x86-64 machine
and the other seven on the arm64 machine. The battery ran on the sources as they stood on 13 October;
since then only comments (docstrings) of \lean{lester} and \lean{kosnita} have changed, with line numbers
preserved. \lean{cliffordchain} was re-run on 14 October at the published commit, including the bridge
and classification files, and the figures below for it are from that run. The gates and tampering
suites were re-run on the final sources on both machines.
\begin{itemize}
\item \emph{Conclusion mutations.} The mutation generator alters conclusions of the forms it supports
(equalities negated, numeric literals changed) and re-checks the proof. It produced 656 mutants; 618
were rejected by an error, 36 by the heartbeat limit (which shows only that the mutated proof did not
complete within the limit), and two survived, both still true: in an auxiliary sign lemma of
\lean{thebault}, the conclusion $\sigma X<0$ weakened to $\sigma X<1$; and in \lean{cliffordchain},
$\mathrm{CliffordChain}(1)$ changed to $\mathrm{CliffordChain}(2)$, which also holds by definition (the
classification proof that uses the lemma then failed to build). Per development:
\lean{clifford} 129/129, \lean{cliffordchain} 283/284, \lean{lester} 82/82, \lean{thebault} 64/65,
\lean{fivecircles} 36/36, \lean{kosnita} 33/33, \lean{blundon} 21/21, \lean{steinerlehmus} 6/6.
\emph{Most main theorems produced no mutants}, because their conclusions (membership, \mdecl{Collinear},
\mdecl{Concyclic}, disjunctions and existentials of these) have no form the generator alters:
\decl{Clifford.clifford_one}, \decl{Clifford.clifford_two}, \decl{Clifford.clifford_two_explicit},
\decl{Clifford.clifford_three}, the second development's \lean{clifford\_one}, \lean{clifford\_two} and
\decl{Clifford.clifford_two_clean}, \decl{Lester.lester}, \decl{Thebault3.sawayama_thebault},
\decl{Fivecircles.fivecircles} and \decl{Kosnita.kosnita}. The mutants that were generated test mainly
the auxiliary algebra, and \decl{Blundon.blundon} and \decl{SteinerLehmus.steiner_lehmus}. (The
\lean{cleanand} tampering test of the second development does check by hand that replacing the
disjunction of \decl{Clifford.clifford_two_clean} by a conjunction breaks the build.)
\item \emph{Contradiction probe.} For each declaration with hypotheses, automation tries to derive
\lean{False} from the hypotheses alone. In every development the corrected probes found no
contradiction. This does not show that the hypotheses are consistent. (A first run of
\lean{clifford} reported 24 contradictions. All were proof checks killed by the operating system for
lack of memory, which the probe had misread as success; 87 mutation checks were killed in the same way
and left unclassified. We changed the harness to count any abnormal exit as an error, removed those
111 records, and re-ran them with less parallelism. The pre-fix result file, which contains the 24 and
the 87 records, is in the verification repository with its SHA-256 checksum.)
\item \emph{Satisfiability.} The battery tries to certify, by a Lean witness, that a declaration's
hypotheses can hold together. Of 362 declarations, 241 were certified: 57 with hypotheses, 150 whose
binders are hypothesis-free, and 34 closed statements (the split is taken from the per-declaration
result files); the other 121 are recorded as unresolved. The closed statements include
\decl{Clifford.cliffordChain_four}, \decl{Clifford.cliffordChain_five} and
\decl{Clifford.Counter6.not_cliffordChain_six}, for which ``certified'' is trivial; the classification
theorems are certified too, which says nothing beyond their being stated. Every main theorem of
Sections~\ref{sec:clifford}--\ref{sec:companions} \emph{with} hypotheses is among the unresolved. So the battery does not
certify that the hypotheses of any main theorem can be satisfied. The numerical re-check above exhibits configurations satisfying all hypotheses of
each main theorem (with the margins stated there); this is evidence, not a certificate. In version~2
the Layer-14 witnesses (Section~\ref{sec:l1415}) are such certificates, in Lean, for every main
theorem except Theorem~\ref{thm:cliff6} and the general-$n$ classification.
\item \emph{Hypothesis ablation.} Each hypothesis \emph{binder} is removed in turn, and we check whether
the unchanged proof still goes through. This tests binders, not their parts: a bundled hypothesis such
as \lean{GenPos5} (many conditions in one binder) or the Fermat-point predicates (three line
memberships each; the Lester proof uses two of the three) counts as one. In \lean{cliffordchain},
\lean{lester}, \lean{thebault}, \lean{fivecircles}, \lean{blundon} and \lean{steinerlehmus} every tested
binder was needed. In \lean{kosnita} two binders of the auxiliary lemma \lean{cert\_den} were not
(they carry a leading underscore for that reason). In \lean{clifford}, \decl{Clifford.clifford_one},
\decl{Clifford.clifford_two} and \decl{Clifford.clifford_three} needed every binder (10, 20 and 41), but \decl{Clifford.clifford_two_explicit} did not need
five of its ten non-collinearity hypotheses (\lean{n135}, \lean{n145}, \lean{n235}, \lean{n245},
\lean{n345}; it has twenty hypotheses in all), the auxiliary lemma \lean{five\_NN} six of its, and \lean{ZD\_ne\_zero} four (all are
non-collinearity hypotheses for triples involving circle 5 or 6; we also confirmed the first two cases
by reading the proofs). These
theorems are stated under more hypotheses than their proofs use; they remain correct. A binder the
unchanged proof needs is not thereby shown to be logically necessary.
\item \emph{Linter.} The unused-variable linter reports no warnings in the library files of all eight
developments as battery-tested (for \lean{cliffordchain}, at the published commit).
\end{itemize}
These checks show that the proofs are accepted and use only the standard axioms, and give evidence
that statements and hypotheses are not vacuous or idle. They do not show that a formal statement means
what the prose says. That question is addressed only partly: by the independent re-verification above
and by the blind re-formalisation of Section~\ref{sec:blind} and the author's
statement-by-statement check.

\paragraph{Statements checked by hand.} The author read each of the fourteen main statements as given in
this paper against its Lean statement at the published commit (a sheet pairing the two, with links to
the source lines, is in the verification repository) and found them to agree. This is a reading check, not a
proof; Section~\ref{sec:blind} explains why the explicit-point theorems must be read with the explicit
point.

\subsection{Layers 14 and 15}\label{sec:l1415}
Version~2 adds two further checks, each in a new file in every repository: \lean{<Lib>/L14.lean} and
\lean{<Lib>/L15Classical.lean}, the latter in namespace \lean{<Lib>.Classical}, for
\lean{<Lib>} $\in\{$\lean{Clifford}, \lean{Fivecircles}, \lean{Lester}, \lean{Thebault3}, \lean{Kosnita},
\lean{Blundon}, \lean{SteinerLehmus}$\}$ (both Clifford repositories use \lean{Clifford}). Their theorems
are added to the gate lists. With them the gates pass, with only the three standard axioms, on
\lean{clifford} 63 declarations, \lean{cliffordchain} 63, \lean{fivecircles} 61, \lean{lester} 75,
\lean{thebault} 34, \lean{kosnita} 25, \lean{blundon} 15 and \lean{steinerlehmus} 10 (346 in all), and
every tampering test is rejected. No pre-existing library file changed except the root file, which gained the two import lines. Nothing was dropped, and no
\lean{set\_option} is used: the two proofs that exceeded the default heartbeat limit were split into
lemmas.

\paragraph{Layer 14: Mathlib-only statements and witnesses.} Each main theorem was restated using only
Mathlib's types and notions, with no constant of our development in its elaborated type (a mechanical
check), and proved from the library theorem through bridge lemmas. This was done for twelve statements:
Theorems~\ref{thm:cliff4}, \ref{thm:cliff5} (one restatement for both Lean forms) and~\ref{thm:cliff6};
$\mathrm{CliffordChain}(4)$, $\mathrm{CliffordChain}(5)$ and the six-circle failure
(\decl{Clifford.l14_not_cliffordChain_six}); and Theorems~\ref{thm:five}--\ref{thm:steiner}. In these, circles
are Mathlib spheres through an arbitrary common point $M$, each $P_{ij}$ is \emph{any} common point other
than $M$, the circles $\Gamma_{ijk}$ are expressed by \mdecl{Cospherical}, the Fermat points are defined by
geometric apexes with strict side conditions (\decl{Lester.l14b_apexOut} identifies them with ours), the
Th\'ebault circles use Mathlib's tangency notions, and circumcentre, incentre, inradius and nine-point
circle are Mathlib's. Eleven of the twelve have a witness theorem that proves every hypothesis on an
explicit configuration together with non-degeneracy of the conclusion (for example, for
Theorem~\ref{thm:cliff5} five pairwise distinct, non-collinear Clifford points; for Blundon both bounds
strict; for Steiner--Lehmus the witness is necessarily isosceles). The exception is
Theorem~\ref{thm:cliff6}: its restatement quantifies the non-collinearity hypotheses over every common
point of the four circles, and discharging them for a concrete configuration would need uniqueness of the
Clifford point, which is not proved. \decl{Clifford.not_cliffordChain_of_six_le} and
\decl{Clifford.cliffordChain_iff} were not restated, since a Mathlib-only form would have to re-encode the
recursive predicate inline; the $n=6$ case, which carries the negative content, is restated. No
restatement was found false. Apart from the Clifford $n=6$ restatement, whose universally quantified
non-collinearity hypotheses are stronger than the library's (equivalent to them given uniqueness of the
Clifford point, which is not formalized), none needed a hypothesis beyond the library theorem's. The Th\'ebault restatement uses Mathlib's
tangency, which implies our distance conditions, so it follows from the library theorem but is not shown
equivalent to it.

\paragraph{Layer 15: textbook statements.} The author collected textbook statements of the seven
results (Clifford's theorems $n=4,5,6$ and the chain, Miquel's five circles, Lester, Sawayama--Th\'ebault,
Kosnita, Blundon, Steiner--Lehmus) from the literature, with their sources. An AI model (an Opus
command-line agent in an isolated folder) was given only these statements and wrote Lean statements for
them; it saw no code and no paper, and the accepted run received the statements in its prompt and otherwise
accessed only its own generated statements, Mathlib probes and compiler output (a transcript audit
in the verification repository found no other file read and no web access). It wrote 18 statements: twelve main statements (Clifford (a), (b), (c), the chain (d) for even
and for odd $n$, the five circles, Lester, Th\'ebault, Kosnita, Blundon's two bounds, Steiner--Lehmus)
and six existence or well-formedness obligations for the constructed objects, marking extra hypotheses
and ambiguities. A comparer then tried to prove each statement, with its exact signature, from the
developments, or else to classify it as covered under extra premises (narrower), out of scope, or false
as written. After a first pass and a gap-closing pass the result is: 16 of the 18 proved exactly (ten
main statements and all six obligations); none narrower; the chain (d) for all $n$ out of scope (only odd
$n=3$ is proved exactly; our $n\le5$ results concern a different predicate); and no writer statement false
as written. New proofs, not in the published developments:
\begin{itemize}
\item Kosnita, the parallel case $AB^2+BC^2+CA^2=5R^2$, which \decl{Kosnita.kosnita} excludes: the lines
are parallel there, so the writer's ``concurrent or all parallel'' holds
(\decl{Kosnita.Classical.cov_Kosnita_zero_locus}, \decl{Kosnita.Classical.cov_D5}).
\item Five circles with no distinctness hypotheses beyond $Q_i\neq V_i$: the distinctness of the other
points follows from the lines being in general position, and injectivity of the five points $Q_i$ is not
needed (\decl{Fivecircles.Classical.gap_D2_distinct}, \decl{Fivecircles.Classical.cov_D2}); the writer
assumes the $Q_i$ not collinear and concludes ``on one circle''.
\item Clifford $n=6$ for the writer's circles $C_m$ through five Clifford points: the identification with
our canonical circles is derived (\decl{Clifford.Classical.gap_canon},
\decl{Clifford.Classical.cov_D1_c}), using the writer's hypothesis that each Clifford point is unique.
\item Lester with the Fermat points defined by equilateral apexes and side conditions
(\decl{Lester.Classical.cov_D3}).
\item Sawayama--Th\'ebault for either assignment of the two circles to the sides of $AD$
(\decl{Thebault3.Classical.cov_D4}).
\item Existence: the pentagram vertices and non-degenerate outer triangles
(\decl{Fivecircles.Classical.cov_D2_obligation}); both Fermat points of a scalene triangle
(\decl{Lester.Classical.cov_D3_obligation}); the two Th\'ebault circles, by an intermediate-value argument
(\decl{Thebault3.Classical.cov_D4_obligation}); $O_a$ when $O$ is off $BC$
(\decl{Kosnita.Classical.cov_D5_obligation}); $R^2-2Rr\ge0$, extracted from the Blundon certificate
(\decl{Blundon.Classical.cov_D6_obligation}); the foot of an angle bisector
(\decl{SteinerLehmus.Classical.cov_D7_obligation}).
\item The chain: for the six centres of Theorem~\ref{thm:chainfail}, every level-$5$ object of the
subfamily omitting circle~$6$ lies on no circle
(\decl{Clifford.Classical.cov_D1_d_no_spherical_level_five}), so the configuration does not meet the
textbook hypotheses (Section~\ref{sec:chainfail}; the identification of the two chain structures is a
premise there).
\end{itemize}
Clifford $n=4,5$, Blundon and Steiner--Lehmus are direct specializations
(\decl{Clifford.Classical.cov_D1_a}, \decl{Clifford.Classical.cov_D1_b}, \decl{Blundon.Classical.cov_D6_lower},
\decl{Blundon.Classical.cov_D6_upper}, \decl{SteinerLehmus.Classical.cov_D7}). For Clifford $n=5$ the
writer additionally assumes that each four-circle common point is unique and that the five points are
not collinear; these assumptions identify the supplied points with ours and remove the line alternative.
Not covered: the chain (d) for
general $n$ in the writer's form, which assumes at every lower level a genuine circle through
non-collinear points or a unique common point. Our $\mathrm{CliffordChain}$ is weaker in its hypotheses,
its failure for $n\ge6$ does not refute the writer's statement, and its $n\le5$ results do not prove it.
The writer's statements are copied verbatim into the Layer-15 files as propositions, and each exact
coverage theorem proves that proposition. Earlier attempts at the writer step by another model were
voided and are not used.

\section{Limitations and further work}\label{sec:limits}
\begin{enumerate}
\item \textbf{Clifford's chain for general $n$} is not proved. The finite-plane predicate we wrote for it
is false for every $n\ge6$ (Theorem~\ref{thm:chainfail}), but that does not refute the textbook chain
with genuine circles and unique points at every level, which remains open here (Section~\ref{sec:l1415});
a corrected statement, on the Riemann sphere
or for $n$ lines, is formulated in prose only. Clifford's theorem itself is proved for $n\le6$, and for
$n=6$ only when the five-point circles are not lines.
\item \textbf{Euclidean, not inversive.} All statements are about Euclidean circles; lines enter only
as stated alternatives. An inversive-plane formulation would remove several exclusions.
\item \textbf{Existence.} The main theorems of Lester and Sawayama--Th\'ebault assume the Fermat points and the Th\'ebault
circles; their existence is proved only in the Layer-15 files, for the definitions used there
(Section~\ref{sec:l1415}).
\item \textbf{Blind checks} establish equivalence of statements, not correctness of the prose; the
blind writer's isolation is shown by its tool log in the verification repository, not enforced; two of the twelve comparisons are
conditional on a fact (uniqueness of the Clifford point) that we have not proved; and
Theorem~\ref{thm:chainfail} was not blind-checked. The Layer-14 and Layer-15 work was likewise done by AI
agents, the Layer-15 writer's isolation is shown by its audited transcript, and no human has checked the
textbook statements against the literature; the Layer-15 Clifford $n=6$ comparison uses the writer's
uniqueness hypothesis, and the chain for general $n$ was not covered.
\end{enumerate}

\section*{Acknowledgements and use of AI}
A large-language-model AI assistant was used to write the Lean proofs, to run the prior-art searches
and the verification battery, and to draft this paper. The author directed the work and reviewed the
paper. Every proof is checked by the Lean kernel, and the author takes responsibility for all content.

\begin{thebibliography}{99}

\bibitem{miquel1838}
A.~Miquel,
\emph{Th\'eor\`emes sur les intersections des cercles et des sph\`eres},
Journal de Math\'ematiques Pures et Appliqu\'ees \textbf{3} (1838), 517--522.

\bibitem{clifford1871}
W.~K.~Clifford,
\emph{Synthetic proof of Miquel's theorem},
The Oxford, Cambridge, and Dublin Messenger of Mathematics \textbf{5} (1871), 124--141.

\bibitem{johnson1929}
R.~A.~Johnson,
\emph{Modern Geometry: An Elementary Treatise on the Geometry of the Triangle and the Circle},
Houghton Mifflin, Boston, 1929.

\bibitem{pedoe1970}
D.~Pedoe,
\emph{A Course of Geometry for Colleges and Universities},
Cambridge University Press, 1970.

\bibitem{lester1997}
J.~A.~Lester,
\emph{Triangles III: Complex triangle functions},
Aequationes Mathematicae \textbf{53} (1997), 4--35.

\bibitem{yiu2010}
P.~Yiu,
\emph{The circles of Lester, Evans, Parry, and their generalizations},
Forum Geometricorum \textbf{10} (2010), 175--209.

\bibitem{sawayama1905}
Y.~Sawayama,
\emph{A new geometrical proposition},
American Mathematical Monthly \textbf{12} (1905), 222--224.

\bibitem{thebault1938}
V.~Th\'ebault,
Problem 3887,
American Mathematical Monthly \textbf{45} (1938), 482--483.

\bibitem{ayme2003}
J.-L.~Ayme,
\emph{Sawayama and Th\'ebault's theorem},
Forum Geometricorum \textbf{3} (2003), 225--229.

\bibitem{grinberg2003}
D.~Grinberg,
\emph{On the Kosnita point and the reflection triangle},
Forum Geometricorum \textbf{3} (2003), 105--111.

\bibitem{etc}
C.~Kimberling,
\emph{Encyclopedia of Triangle Centers},
\url{https://faculty.evansville.edu/ck6/encyclopedia/ETC.html}.

\bibitem{blundon1965}
W.~J.~Blundon,
\emph{Inequalities associated with the triangle},
Canadian Mathematical Bulletin \textbf{8} (1965), 615--626.

\bibitem{gilbert1963}
G.~Gilbert and D.~MacDonnell,
\emph{The Steiner--Lehmus theorem},
American Mathematical Monthly \textbf{70} (1963), 79--80.

\bibitem{coxeter1967}
H.~S.~M.~Coxeter and S.~L.~Greitzer,
\emph{Geometry Revisited},
Mathematical Association of America, 1967.

\bibitem{chou1988}
S.-C.~Chou,
\emph{Mechanical Geometry Theorem Proving},
D.~Reidel, Dordrecht, 1988.

\bibitem{afpmiquel}
A.~Freitas Ramos, D.~Barros Hulak and R.~J.~G.~B.~de Queiroz,
\emph{Miquel's Theorem in Isabelle/HOL}, Archive of Formal Proofs, 11 August 2026.
\url{https://www.isa-afp.org/entries/Miquel.html}

\bibitem{hsgeometry}
\emph{HighSchoolGeometry} (Rocq library), file \texttt{applications\_cocyclicite.v}.
\url{https://github.com/rocq-community/HighSchoolGeometry}

\bibitem{leangenius}
R.~Walters,
\emph{lean-genius}, GitHub repository, pull request 25056, 2026.
\url{https://github.com/rjwalters/lean-genius/pull/25056}

\bibitem{leangeniusSL}
R.~Walters,
\emph{lean-genius}, file \texttt{proofs/Proofs/SteinerLehmusTheoremOQ01.lean}, pull request 25041, 2026.
\url{https://github.com/rjwalters/lean-genius/blob/main/proofs/Proofs/SteinerLehmusTheoremOQ01.lean}

\bibitem{leangeniusG}
R.~Walters,
\emph{lean-genius}, file \texttt{proofs/Proofs/HeronsFormulaOQ06OQ02.lean} (Gerretsen's inequality), 2026.
\url{https://github.com/rjwalters/lean-genius/blob/main/proofs/Proofs/HeronsFormulaOQ06OQ02.lean}

\bibitem{acornSL}
\emph{acornlib} (Acorn prover library), file \texttt{src/theorems1000/theorem\_steiner\_lehmus.ac}.
\url{https://github.com/acornprover/acornlib/blob/master/src/theorems1000/theorem_steiner_lehmus.ac}

\bibitem{mathlib1000}
The mathlib Community, \emph{1000+ theorems in Lean}, \texttt{docs/1000.yaml}.
\url{https://leanprover-community.github.io/1000.html}

\bibitem{afp}
\emph{Archive of Formal Proofs}, ISSN 2150-914x. \url{https://isa-afp.org}

\bibitem{lean4}
L.~de Moura and S.~Ullrich,
\emph{The Lean 4 theorem prover and programming language},
in: Automated Deduction -- CADE-28, Lecture Notes in Computer Science \textbf{12699}, Springer, 2021, 625--635.

\bibitem{mathlib2020}
The mathlib Community,
\emph{The Lean mathematical library},
in: Proceedings of the 9th ACM SIGPLAN International Conference on Certified Programs and Proofs (CPP '20), ACM, 2020, 367--381.

\bibitem{paperA}
J.~Bald,
\emph{Classical plane geometry by coordinate certificates: Desargues, Monge, Simson, Napoleon and Lami in Lean~4},
Zenodo, 2026. \url{https://doi.org/10.5281/zenodo.23239930}

\bibitem{paperB}
J.~Bald,
\emph{The Malfatti circles in Lean~4, with independent certificate proofs of Feuerbach, Miquel and Brahmagupta},
Zenodo, 2026. \url{https://doi.org/10.5281/zenodo.23240500}

\bibitem{paperC}
J.~Bald,
\emph{Barbier's theorem for the Reuleaux triangle and the convex four-vertex theorem in Lean~4, through support functions},
Zenodo, 2026. \url{https://doi.org/10.5281/zenodo.23323062}

\end{thebibliography}

\end{document}
